% !TEX root = ../main.tex

\chapter{参考值的计算与可靠性说明}\label{ch:reference_values}
  为了更客观地评估第\ref{ch:numerical}章中的数值结果, 本附录说明各算例参考值的来源.
  对具有显式解的方程, 直接使用解析公式;
  对没有简单闭式解的方程, 则采用独立于正文求解器的数值方法构造高精度基准,
  并同时给出误差或置信区间估计.

  \section{Heat 方程的显式参考值}\label{sec:ref_heat}
    对 Heat 方程,
    \begin{equation}
      g(x)=\exp\!\left(-\frac{\lVert x\rVert^2}{d}\right),
      \qquad x_0=0,\qquad T=1.
    \end{equation}
    由正文中的 \eqref{eq:numerical_heat_solution} 可得
    \begin{equation}
      Y_0 = \left(\frac{d}{d+2}\right)^{d/2}.
    \end{equation}
    因而
    \begin{equation}
      d=20 \Rightarrow Y_0=0.385543289429532,
    \end{equation}
    \begin{equation}
      d=50 \Rightarrow Y_0=0.375116802253964,
    \end{equation}
    \begin{equation}
      d=100 \Rightarrow Y_0=0.371527882126961.
    \end{equation}
    这些值由显式公式直接给出, 不涉及额外数值误差.

  \section{HJB-LQ 方程的积分参考值}\label{sec:ref_hjb}
    对 HJB-LQ 方程, 由正文中的 \eqref{eq:numerical_hjblq_reference},
    \begin{equation}
      Y_0
      =
      -\log
      \mathbb{E}
      \left[
        \frac{2}{1+2T\chi_d^2}
      \right].
    \end{equation}
    因此问题可以化为单变量积分
    \begin{equation}
      \mathbb{E}
      \left[
        \frac{2}{1+2T\chi_d^2}
      \right]
      =
      \int_{0}^{\infty}
      \frac{2}{1+2Ty} f_{\chi_d^2}(y)\,\dif y.
    \end{equation}
    本文采用自适应 Simpson 积分在有限区间上数值计算该积分.
    当积分上限从 $300$ 增加到 $500$ 时, 结果变化不超过 $5\times 10^{-12}$,
    因而截断误差可以忽略.

    最终得到
    \begin{equation}
      d=50 \Rightarrow Y_0=3.882006682195226,
    \end{equation}
    \begin{equation}
      d=100 \Rightarrow Y_0=4.590161724604434.
    \end{equation}

  \section{Black--Scholes 方程的 Monte Carlo 参考值}\label{sec:ref_bsm}
    对 Black--Scholes 方程, 终端 payoff 为高维算术平均型欧式看涨期权
    \begin{equation}
      g(x)=\max\!\left(\frac{1}{d}\sum_{m=1}^{d}x_m-K,0\right).
    \end{equation}
    该问题没有与一维情形对应的简单闭式解.
    因此本文采用独立的大样本 Monte Carlo 估计其初值,
    并使用几何平均期权的闭式价格作为控制变量.

    具体地说, 记算术平均 payoff 对应的贴现随机变量为 $A$,
    几何平均 payoff 对应的贴现随机变量为 $G$.
    由于几何平均
    \begin{equation}
      \bar{G}_T = \exp\!\left(\frac{1}{d}\sum_{m=1}^{d}\log X_T^{m}\right)
    \end{equation}
    仍服从对数正态分布, 因而 $\mathbb{E}[G]$ 可以显式计算.
    于是采用控制变量估计量
    \begin{equation}
      A^{\ast}=A-\beta\bigl(G-\mathbb{E}[G]\bigr),
    \end{equation}
    同时结合对偶路径进一步降低方差.

    本文使用 $10^6$ 条路径计算参考值, 并报告 $95\%$ 置信区间:
    \begin{equation}
      d=50 \Rightarrow Y_0=0.04921354,\qquad
      95\%\text{ CI}=[0.04920245,\ 0.04922462],
    \end{equation}
    \begin{equation}
      d=100 \Rightarrow Y_0=0.04880920,\qquad
      95\%\text{ CI}=[0.04880137,\ 0.04881702].
    \end{equation}

  \section{Allen--Cahn 方程的有限差分参考值}\label{sec:ref_allencahn}
    对 Allen--Cahn 方程, 本文不采用正文中的随机特征求解器来构造参考值,
    而是利用其终端条件与扩散过程均具有径向对称性这一事实,
    将问题转化为径向变量上的一维半线性抛物方程.
    记 $r=\lVert x\rVert$, 则解可写为 $u(t,x)=v(t,r)$.
    在反向时间变量 $\tau=T-t$ 下, 参考解满足
    \begin{equation}
      \frac{\partial v}{\partial \tau}
      =
      \frac{\partial^2 v}{\partial r^2}
      + \frac{d-1}{r}\frac{\partial v}{\partial r}
      + v - v^3,
      \qquad
      v(0,r)=\frac{1}{2+0.4r^2}.
    \end{equation}
    本文在区间 $r\in[0,20]$ 上采用径向有限差分离散,
    对扩散项作隐式处理, 对非线性反应项作显式处理.
    在 $r=0$ 处施加径向对称边界条件, 在远端边界取零值近似.

    为评估离散误差, 分别使用两组网格
    \begin{equation}
      (\Delta r,\Delta \tau)=(0.025,\ 5\times 10^{-5}),
    \end{equation}
    \begin{equation}
      (\Delta r,\Delta \tau)=(0.0125,\ 2.5\times 10^{-5})
    \end{equation}
    进行计算, 并将两组结果的差作为误差估计.
    最终采用细网格结果作为参考值, 得到
    \begin{equation}
      d=50 \Rightarrow Y_0=0.09908593,\qquad
      \text{网格差约 }3.19\times 10^{-6},
    \end{equation}
    \begin{equation}
      d=100 \Rightarrow Y_0=0.05278464,\qquad
      \text{网格差约 }2.81\times 10^{-6}.
    \end{equation}
    此外, 当边界截断半径从 $15$ 增加到 $30$ 时结果几乎不变,
    因而边界截断误差相较于网格误差可以忽略.
